% The proposed runtime. Sources: space_runtime/{agent,facts,context,guard,verification,report}.py, prompts/ours/
% {system,step}.md and tools.yaml (F124), openral_ext/capability_server/contracts.yaml, configs/space_runtime.yaml,
% experiments/operator_console/console.py (F125), space_context/link.py; framing from AGENT/research_content.md §4 and
% the user's direction (F148: trustworthy AI, the harness contains the consequence). Implementation names (simulator,
% middleware, planners, model names) are kept out of this section; they are in Section setup.
\providecommand{\TBD}[1]{\textcolor{red}{[TBD: #1]}}

\section{The proposed runtime}
\label{sec:method}

The proposed runtime puts a foundation model where a mission executive would run an authored procedure. Every
responsibility that space operations assign to independently engineered mechanisms stays outside the model. The model
decides what to do next. It chooses what to observe, how to break the objective into tasks and which capability to
call, and it decides how to recover and when to involve mission control. The runtime decides what is true, what is
permitted and what physically happened. This division meets R1 by construction, since no part of the runtime chooses an action. It also fixes where
trust comes from. A model can misread a scene, lose track of an earlier result or propose a reserved operation. The
design question is which mechanism contains the consequence when that happens. In the proposed runtime that mechanism is
never the model itself (Fig.~\ref{fig:architecture}).

\begin{figure}[t]
  \centering
  \fbox{\parbox[c][4cm][c]{0.9\columnwidth}{\centering architecture figure placeholder}}
  \caption{\TBD{The step loop of the proposed runtime: mission record, step context, model, guard, capability boundary,
  verifier, trace; mission control behind the link. Mark which maps are outside the model.}}
  \label{fig:architecture}
\end{figure}

One step of the runtime is a fixed chain of maps. The runtime keeps an append-only record $E_t=(e_1,\dots,e_t)$ of
typed events. A reducer $\phi$ folds the record into the mission facts $F_t=\phi(E_t)$. A renderer $\rho$ turns the
facts into the step context $c_t=\rho(F_t)$, which is everything the model reads. The model proposes exactly one call
$a_t$ from a closed set $\mathcal{A}$ of typed tools. When $a_t$ acts on the robot, the guard returns an outcome
$g(a_t,F_t)\in\{\text{allow},\text{deny},\text{approval required},\text{unavailable}\}$. Only an allowed call reaches
the capability boundary, the same boundary the baselines call. The boundary executes the call and returns
robot-observable evidence $y_t$. The verifier maps the call and its evidence to effect records $v(a_t,y_t)$. Each record
has a verdict that holds, fails or is not established. Every outcome is appended to the record, and the next step
starts again from $\phi$. The model's output enters the record only as a proposal. Four maps stand between that proposal
and mission truth, namely $\phi$, $\rho$, $g$ and $v$. Together they form the harness around the model, and each answers
a requirement of Section~\ref{sec:requirements}.

The reducer $\phi$ makes the record the only mission truth (R2). Facts change only through typed events. These are the
objective as received, execution results with their evidence, verified effects and answers to queries. They also
include guard decisions, fault and availability changes, and mission control's answers and approvals. The state of a
manipulated item changes only through a verified effect, and it carries the effect that set it. No claim of the model
moves an item. The model's working plan is kept in the record as working state. The model writes and revises it, and
nothing reads it as evidence. Truth lives in the record and not in a conversation. A reasoning service that stops and
restarts therefore resumes from the same facts, and so does a different service that takes over. A call that was
running when the previous reasoner stopped becomes an anomaly of its own. The successor sees the gap instead of assuming
the result.

The renderer $\rho$ rebuilds the step context from the facts at every step, and the model has no memory beyond it. A
standing text is fixed for each reasoner session. It gives the operating rules and the authority model, which is read
off the capability contracts and the guard configuration. It also gives the constants of the link to Earth, the
surveyed site and the equipment catalogue. The step message gives the record. It holds the objective, the robot's latest
evidence, the verified effects, the observations and the state of each manipulated item. It also holds the manipulation
attempts, the unresolved anomalies, the exchanges with mission control, the state of the link, the active faults, the
working plan and the last 12 calls. The renderer selects and formats facts. It writes no advice and names no situation.
Two properties follow. A restarted or replaced reasoner reads exactly what its predecessor would have read. And the
context grows with the number of items and open anomalies, not with the number of steps, because an old result leaves
the recent window while its effect stays in the facts.
% "last 12 calls": space_runtime/facts.py RECENT = 12. Sections of the step message: prompts/ours/step.md; windows in
% space_runtime/context.py (verified effects 12, observations 8, attempts 10; items, labels, anomalies in full).

The model acts only through the closed set $\mathcal{A}$ of whole operations with mission-level arguments. It can drive
to a surveyed feature or a pose and dock at a feature. It can ready or stow the arm, pick an item by its catalogued
handle, place an item on a surveyed support, press a declared button and open the jaws. It observes through read-only
queries of the site survey, the serviced units' own telemetry, mission control's inventory record and the robot's
state. A further query asks a vision-language model a question about one camera frame. The model reads text only. An
image reaches it as the vision model's answer to a question the model asked. That answer enters the record as an
observation, never as a fact about the world. Each capability has a typed contract that says what one call means and
what it returns, with no ordering and no recovery. The mechanics of an operation live in the skills below the contract,
which every method shares. Examples are where the camera looks, how a grasp is fitted and how the arm withdraws. The
model therefore chooses among operations. It cannot command a joint, create a capability or write a fact.
% Mission-level effects and queries: contracts.yaml entries without `engineering: true`; the local tools (working plan,
% report, request, wait) are prompts/ours/tools.yaml. No image is sent to the reasoner (space_runtime/tools.py).

The guard $g$ enforces authority and invariant boundaries outside the model (R4). Before an effect is dispatched, it
checks explicit boundaries only. The capability must exist, and its arguments must satisfy the contract. It must not be
unavailable, no active fault may prohibit it, and no configured hard invariant may be violated. Its authority must
permit it. Authority has two sources. Each capability carries an authority class, and a target rule classes a call by
what it acts on. The rule exercised in this study reserves the removal of a module from the bay of a unit that reports
itself energized. It follows the flight-rule principle that equipment is not removed or altered without authorization
\citep{cfr1214403}. The guard reads the units' telemetry fresh before it decides. A held call is recorded as pending. A
request to mission control sent while the call is pending carries it as structured data. A grant is recorded as an
approval of that capability on that target, and the guard then allows the call. The outcome and its reason return to
the model as context. The guard never judges whether a call is a good next step, so it cannot become a second planner.
Its value is a guarantee that does not depend on the model. A model that has misunderstood the rule still cannot
execute the removal.

The verifier $v$ establishes what physically happened (R3). It compares the evidence the boundary returns with the goal
the call commanded. It never reads the skill's success flag, the model's claims, the order of calls or simulator truth.
A drive holds when the rover's own localization estimate lies within the navigation tolerance of the goal. A pick holds
when the jaw encoder, the arm kinematics and the wrist images before and after the lift agree that the item rose with
the gripper. A place holds when the jaws are open, the tool is clear, the item stays still in the settle frames and no
depth points remain between the fingers. A press holds when the tool went past the fitted button's tip and withdrew.
The verifier recomputes the image checks from the raw frames the skill recorded, so the skill's own verdict is never its
evidence. Its verdict has three values. An effect that the evidence neither shows nor contradicts is recorded as not
established. The item's state then becomes unknown, which leaves the model to look. Whatever the record says is carried
is examined again after every call that does not itself let go. The vision model reads the wrist frame, and the jaw
encoder can overrule it when the fingers are shut on nothing. A payload can leave the jaws while the rover drives, and
this check catches the loss before an unrelated failure hides it. The model's statement that the mission succeeded is
never evidence. The report to mission control is generated from the record, with the model's own words quoted beside
the verified effects.

Detection and escalation follow the fault-management discipline of onboard autonomy (R5). Failures, refusals and
disagreements between sources become typed anomalies in the record. One such disagreement is a catalogued handle that
the wrist depth does not match. Each anomaly stays until a later success of the same capability resolves it. A latched
fault of the robot's safety kernel prohibits every effect. Any other fault or degradation is a fact for the model to
reason about. The runtime does not decide when to escalate. It gives the model the means and the facts. A request to
mission control declares one of four kinds, an authorization, a check of what lies within the robot's authority, a new
procedure, or information. It carries the calls the guard is holding. The request does not block the robot, because its
answer crosses a delayed link and may be queued until the next pass. The model continues other work and reads the
answer when it arrives. The context also states the link's state, the time until Earth sets and mission control's
decision time. These are facts computed from the site's geometry and the published protocol, and none of them is a
rule. The runtime thus keeps detection, isolation into a typed anomaly and escalation to the next level of authority
\citep{ecss2008operability}. The recovery strategy stays with the model.
% Four request kinds: experiments/operator_console/console.py REQUEST_KINDS; prompts/ours/tools.yaml request_operator.
% Anomaly kinds and resolution: space_runtime/facts.py apply() (a later success of the same capability resolves them).

The record is also the trace (R6). For each call it joins the proposal with the model's one-sentence rationale, the
guard's outcome, the execution result with its evidence, and the verification. Mission control can reconstruct from it
what the robot knew at each step, what it attempted, which boundary acted and what physically happened. It needs no
access to the model's private reasoning for that.
% research_content.md §4.13 (trace chain). Generated report: space_runtime/report.py.

Nothing in the runtime names a situation, a task order or a recovery. It pre-engineers boundaries and evidence, and two
conditions bound what that can deliver. The capability set must contain every physical operation a solution needs,
because the model composes operations and cannot invent one. And each mechanism contains a consequence only where it
has evidence or a rule for it. The verifier can refuse a false effect only when the robot's sensors can establish that
effect. The guard can hold only a call that some rule classes. Removing a mechanism should therefore matter where
nothing else in the system observes or bounds the consequence, and nowhere else. Section~\ref{sec:setup-ablation}
defines the test of this prediction for each of $\phi$, $\rho$, $g$ and $v$.

% Disagreements between this section (as built) and AGENT/research_content.md §4:
% - §4.10 lists "deterministic simulator truth" among the verifier's evidence sources. As built, the verifier never reads
%   simulator truth (space_runtime/verification.py docstring; configs/space_runtime.yaml verification comment); truth is
%   the external evaluator's only.
% - §4.8/§4.9 describe communication, compute, energy, thermal and illumination facts and onboard / edge / Earth reasoning
%   tiers with reconciliation. As built, the space context is the Direct-with-Earth link (SPICE), mission control's
%   decision time, and the units' telemetry (load, criticality, energized, breaker); one remote reasoner; reconstruction
%   after restart is supported (facts.py reasoner_history, agent.py resume option) but no endpoint change occurred in
%   the experiments (tested offline by ablation J3).
% - §4.7 lists hard resource/health limits and staleness checks among guard inputs. As built, `invariants: []` and
%   `fdir_prohibits: {estop: [all_effects]}` (configs/space_runtime.yaml); the only authority rule is the target rule
%   energized_equipment_removal; every capability contract has authority_class routine. The mechanism exists; only the
%   target rule was exercised. Capability unavailability is a guard outcome, but no event in the runs made a capability
%   unavailable (agent.py emits no availability_changed event).
% - §1.4/§2.5 motivate the design from coding-agent harnesses (Hermes, Pi, Codex). Per the user's framing (F148), this
%   section derives the design from space-operations disciplines instead and makes no coding-agent analogy.
